% !TEX root = ../../../main.tex

\subsection{候选序列生成与 DeepSeek 语义修正}

在得到原始 Gloss 序列 \texttt{rawSequence} 后，
系统还需要将词级识别结果转换为用户能够理解的自然语言表达。

由于句子视频识别结果来自滑动窗口分类和词段级后处理，
\texttt{rawSequence} 中仍然可能存在边界插入词、局部词汇误判或表达不够自然等问题。

因此，
本系统在直接中文映射的基础上，
引入 DeepSeek 语义修正模块，
用于对后端生成的候选 Gloss 序列进行排序选择，
并将所选序列转换为忠实、自然的英文和中文表达。

需要说明的是，
DeepSeek 并不直接参与视频识别，
也不会重新分析原始图像帧。

视觉识别仍由 Python 端完成，
包括视频帧读取、MediaPipe 关键点检测、词级分类模型推理和词段后处理。

DeepSeek 接收的是后端整理后的文本序列、词段候选信息和程序生成的有限候选序列。

也就是说，
DeepSeek 的作用不是自由改写识别结果，
而是在受到候选序列约束的前提下，
选择更合理的候选序列并完成忠实自然化表达。

\WhuAutoFigure{0.82\textwidth}{0.42\textheight}{figures/chapter06/deepseek-candidate-correction-flow.png}{DeepSeek 候选序列语义修正流程图}{fig:chapter06-deepseek-candidate-correction}

如图~\ref{fig:chapter06-deepseek-candidate-correction} 所示，
语义修正流程可以分为四个阶段。

第一，
Python 识别服务返回原始 Gloss 序列、中文逐词映射结果和每个词段的 Top-K 候选信息。

第二，
后端根据原始词段和 Top-K 候选生成一组合法候选序列。

第三，
后端将原始结果、词段证据和候选序列封装为 DeepSeek 请求。

第四，
DeepSeek 从候选集合中选择一个候选编号，
后端再根据该候选编号生成最终修正结果。

表~\ref{tab:chapter06-deepseek-input-fields} 给出了 DeepSeek 语义修正请求中的主要输入字段、数据来源及其作用。

\begin{longtable}{L{0.20\textwidth} L{0.24\textwidth} L{0.28\textwidth} L{0.20\textwidth}}
  \caption{DeepSeek 语义修正输入信息说明表}
  \label{tab:chapter06-deepseek-input-fields}\\
  \toprule
  输入字段 & 数据来源 & 主要内容 & 作用 \\
  \midrule
  rawSequence & Python 句子视频识别结果 & 原始 Gloss 序列 & 表示视觉识别模块输出的词级识别结果 \\
  rawTextZh & 后端中文映射结果 & 逐词映射后的中文文本 & 作为直接展示结果和自然化翻译参照 \\
  segments & Python 返回的 segmentTopK 经后端整理得到 & 词段编号、原始标签、中文标签、置信度和 Top-K 候选 & 为候选选择提供视觉识别证据 \\
  candidates & 后端候选生成模块 & 由保留、删除和替换操作组合生成的合法候选序列 & 约束 DeepSeek 只能从程序候选中选择，避免自由改写 \\
  \bottomrule
\end{longtable}

DeepSeek 请求中的 \texttt{segments} 字段来源于 Python 侧返回的 \texttt{segmentTopK}。

每个词段不仅包含原始识别标签，
还包含该词段范围内的 Top-K 候选标签及其概率统计。

这些候选信息为后端生成替换候选提供依据，
也为 DeepSeek 判断某个替换是否具有视觉证据提供参考。

在字段结构上，
语义修正请求可以抽象为代码清单~\ref{lst:chapter06-deepseek-input-example} 所示形式。

\begin{lstlisting}[style=whuCode,caption={DeepSeek 语义修正输入结构示例},label={lst:chapter06-deepseek-input-example}]
{
  "rawSequence": ["you", "want", "learn"],
  "rawTextZh": "你 想要 学习",
  "segments": [
    {
      "segmentIndex": 1,
      "rawLabel": "you",
      "rawLabelZh": "你",
      "avgConfidence": 0.86,
      "maxConfidence": 0.94,
      "topK": [
        {
          "label": "you",
          "labelZh": "你",
          "avgProb": 0.86
        },
        {
          "label": "who",
          "labelZh": "谁",
          "avgProb": 0.21
        }
      ]
    }
  ],
  "candidates": [
    {
      "candidateId": "C000",
      "sequence": ["you", "want", "learn"],
      "edits": [],
      "editCount": 0,
      "replaceCount": 0,
      "removeCount": 0,
      "source": "raw"
    },
    {
      "candidateId": "C001",
      "sequence": ["you", "learn"],
      "edits": [
        {
          "type": "remove",
          "segmentIndex": 2,
          "from": "want",
          "to": null
        }
      ],
      "editCount": 1,
      "replaceCount": 0,
      "removeCount": 1,
      "source": "remove"
    }
  ]
}
\end{lstlisting}

后端首先基于原始词段和每个词段的 Top-K 候选构造候选 Gloss 序列。

对于每个词段，
程序会生成三类操作选项。

第一类是保留原标签，
即该词段继续使用视觉识别模块给出的 \texttt{rawLabel}。

第二类是删除该词段，
用于处理明显的边界插入词、重复词或过渡噪声。

第三类是替换该词段，
即使用同一词段 Top-K 候选中的其他标签替代原标签。

随后，
系统对各词段的操作选项进行组合，
形成候选序列集合 \texttt{candidates}。

为了避免候选空间过大，
并防止语义修正幅度过强，
后端对候选生成过程设置了约束。

当前实现中，
总编辑次数最多为 2 次，
替换次数最多为 1 次，
候选序列数量最多为 81 个。

如果候选序列为空，
或者原始序列本身为空，
系统会直接进入 fallback 分支，
不再调用 DeepSeek。

DeepSeek 请求由后端统一构造。

系统提示词将 DeepSeek 约束为手语 Gloss 候选排序器和忠实自然化翻译器，
要求其只能从 \texttt{candidates} 中选择一个 \texttt{candidateId}。

DeepSeek 不允许新增词段，
不允许重排词段，
也不允许使用候选集合之外的 Gloss 词。

同时，
DeepSeek 返回的 \texttt{correctedGlossSequence} 必须与所选候选序列完全一致。

表~\ref{tab:chapter06-deepseek-output-fields} 总结了 DeepSeek 返回结果中的主要字段含义及后端处理方式。

\begin{longtable}{L{0.24\textwidth} L{0.32\textwidth} L{0.34\textwidth}}
  \caption{DeepSeek 语义修正返回字段说明表}
  \label{tab:chapter06-deepseek-output-fields}\\
  \toprule
  返回字段 & 含义 & 系统处理方式 \\
  \midrule
  selectedCandidateId & DeepSeek 从候选集合中选择的候选编号 & 后端校验该编号是否存在于 candidates 中，并以该编号对应的程序候选序列为准 \\
  correctedGlossSequence & DeepSeek 写出的修正后 Gloss 序列 & 理论上应与 selectedCandidateId 对应候选一致；若不一致，后端仍信任 selectedCandidateId 对应候选 \\
  englishSentence & 忠实英文句子或片段 & 用于中间语义解释和调试展示 \\
  chineseTranslation & 忠实自然中文表达 & 作为 correctedTextZh 的主要来源，用于前端展示 \\
  translationConfidence & 翻译置信等级 & 用于提示用户当前自然化结果的可靠程度 \\
  isIncomplete & 语义是否不完整 & 若为 true，说明结果更适合作为片段展示 \\
  translationNote & 不确定性或缺失信息说明 & 用于解释低置信度或语义不完整情况 \\
  reason & 选择候选的简要原因 & 用于调试和分析候选选择依据 \\
  \bottomrule
\end{longtable}

一次正常的 DeepSeek 语义修正返回结果可以抽象为代码清单~\ref{lst:chapter06-deepseek-output-example} 所示形式。

\begin{lstlisting}[style=whuCode,caption={DeepSeek 语义修正返回结构示例},label={lst:chapter06-deepseek-output-example}]
{
  "selectedCandidateId": "C000",
  "correctedGlossSequence": ["you", "want", "learn"],
  "englishSentence": "You want to learn.",
  "chineseTranslation": "你想学习。",
  "translationConfidence": "high",
  "isIncomplete": false,
  "translationNote": "The gloss relation is clear.",
  "reason": "The raw candidate is coherent and supported by segment evidence."
}
\end{lstlisting}

后端接收到 DeepSeek 返回内容后，
不会无条件采用模型写出的 \texttt{correctedGlossSequence}。

系统首先解析 DeepSeek 返回的 JSON，
并读取其中的 \texttt{selectedCandidateId}。

随后，
后端在程序生成的候选集合中查找该编号对应的候选序列。

如果该编号存在，
最终修正后的 Gloss 序列以该候选序列为准。

如果 DeepSeek 写出的 \texttt{correctedGlossSequence} 与所选候选序列不一致，
后端也不会直接信任模型写出的序列，
而是继续使用 \texttt{selectedCandidateId} 对应的程序候选。

这一设计保证了语义修正结果始终被限制在后端生成的候选空间中，
避免大语言模型生成不受视觉证据支持的新词、
新增词段或重排词序。

在候选选择完成后，
后端会根据所选候选生成 \texttt{correctedGlossSequence}，
并结合 DeepSeek 返回的中文表达生成 \texttt{correctedTextZh}。

同时，
系统还会记录 \texttt{selectedSegments} 和 \texttt{removedSegments} 等辅助信息，
用于说明哪些词段被保留、替换或删除。

例如，
当某个候选序列表示删除一个边界插入噪声词时，
DeepSeek 只需要选择对应的候选编号，
后端即可根据该候选的编辑记录生成删除词段说明。

如果 DeepSeek 调用失败、
返回内容无法解析、
缺少 \texttt{selectedCandidateId}，
或返回的候选编号不存在于候选集合中，
系统会进入 fallback 分支。

fallback 分支会保留原始 \texttt{rawSequence} 和 \texttt{rawTextZh}，
并将 \texttt{selectedCandidateId} 设置为 \texttt{C000}。

此时，
系统不会标记语义修正已生效，
而是将结果视为原始识别结果的直接展示。

通过上述机制，
本系统将大语言模型限定在候选排序和忠实自然化表达范围内。

这种设计既利用了 DeepSeek 对短句语义合理性的判断能力，
又避免了大语言模型自由生成导致的识别结果失真。

因此，
候选序列生成与 DeepSeek 语义修正模块并不是替代视觉识别模块，
而是在视觉识别结果基础上，
对有限候选进行保守选择和自然语言表达转换，
从而提升句子视频识别结果的可读性与用户可理解性。